\nomenclature[00]{$\mathbb{R}$}{множество вещественных чисел}
\nomenclature[01]{$\mathbb{X} \subset \mathbb{R}^{n}$}{пространство исходной переменной}
\nomenclature[02]{$\mathbf{x} \in \mathbb{X}$}{вектор признакового описания}
\nomenclature[03]{$\mathbb{Y} \subset \mathbb{R}^{K}$}{пространство целевой переменной}
\nomenclature[04]{$\mathbf{y} \in \mathbb{Y}$}{целевая переменная (вектор)}
\nomenclature[05]{$y \in \mathbb{Y}$}{целевая переменная (скаляр)}
\nomenclature[06]{$p(\mathbf{x}, \mathbf{y})$}{неизвестное распределение данных}
\nomenclature[07]{$\mathfrak{D}_{m} = \left\{ \left( \mathbf{x}_{i}, \mathbf{y}_{i} \right) \right\}_{i=1}^{m}$}{выборка из~$m$ независимых реализаций пар~$\left( \mathbf{x}, \mathbf{y} \right)$ согласно неизвестному распределению~$p(\mathbf{x}, \mathbf{y})$}
\nomenclature[08]{$\mathbb{W} \subset \mathbb{R}^{N}$}{пространство параметров модели}
\nomenclature[09]{$\mathbf{w} \in \mathbb{W}$}{вектор параметров модели}
\nomenclature[10]{$f_{\mathbf{w}}: \mathbb{X} \to \mathbb{Y}$}{параметрическая модель}
\nomenclature[11]{$\ell: \mathbb{Y} \times \mathbb{Y} \to \mathbb{R}$}{функция ошибки}
\nomenclature[12]{$\mathcal{L}_m: \mathbb{W} \to \mathbb{R}$}{эмпирическая функция потерь на выборке~$\mathfrak{D}_{m}$}
\nomenclature[13]{$\mathcal{L}: \mathbb{W} \to \mathbb{R}$}{популяционная функция потерь}
\nomenclature[14]{$p(\mathbf{y} | \mathbf{x}, \mathbf{w})$}{правдоподобие объекта}
\nomenclature[15]{$p(\mathbf{w})$}{априорное распределение параметров модели}
\nomenclature[16]{$p(\mathbf{y}, \mathbf{w} | \mathbf{x})$}{вероятностная параметрическая модель}
\nomenclature[17]{$p(\mathbf{w} | \mathfrak{D}_{m})$}{апостериорное распределение параметров модели}
\nomenclature[18]{$\mathbf{g}_{i}(\mathbf{w}) \in \mathbb{R}^{N}$}{вектор градиента эмпирической функции потерь на~$i$-м объекте}
\nomenclature[19]{$\mathbf{g}^{(m)}(\mathbf{w}) \in \mathbb{R}^{N}$}{вектор градиента эмпирической функции потерь на выборке~$\mathfrak{D}_{m}$}
\nomenclature[20]{$\mathbf{H}_{i}(\mathbf{w}) \in \mathbb{R}^{N \times N}$}{матрица Гессе эмпирической функции потерь на~$i$-м объекте}
\nomenclature[21]{$\mathbf{H}^{(m)}(\mathbf{w}) \in \mathbb{R}^{N \times N}$}{матрица Гессе эмпирической функции потерь на выборке~$\mathfrak{D}_{m}$}
\nomenclature[22]{$\mathbf{z} = f_{\mathbf{w}}(\mathbf{x}) \in \mathbb{R}^{K}$}{выход модели}
\nomenclature[23]{$\mathbf{J}(\mathbf{w}) \in \mathbb{R}^{K \times N}$}{якобиан выхода модели по параметрам}
\nomenclature[24]{$\mathbf{A}(\mathbf{w}) \in \mathbb{R}^{K \times K}$}{матрица кривизны функции потерь по выходу модели}
\nomenclature[25]{$\mathbf{G}(\mathbf{w}) \in \mathbb{R}^{N \times N}$}{гаусс-ньютоновская компонента матрицы Гессе}
\nomenclature[26]{$\mathbf{R}(\mathbf{w}) \in \mathbb{R}^{N \times N}$}{остаточная компонента матрицы Гессе}
\nomenclature[27]{$\mathbf{S}^{(p)}(\mathbf{w}) \in \mathbb{R}^{K \times n_p}$}{матрица распространения возмущения предактивации $p$-го слоя к выходу сети}
\nomenclature[28]{$\mathbf{P}^{(p)}(\mathbf{w}) \in \mathbb{R}^{n_p \times N_p}$}{параметрический якобиан $p$-го слоя}
\nomenclature[29]{$\mathbf{D}^{(p)}(\mathbf{w}) \in \mathbb{R}^{n_p \times n_p}$}{диагональная матрица производных активации}
\nomenclature[30]{$q(\mathbf{w})$}{функция предпочтения по выбору точек для оценки сходимости ландшафта функции потерь}
\nomenclature[31]{$\Delta_{p}(k)$}{критерий $p$-сходимости ландшафта функции потерь}
\nomenclature[32]{$\Delta_{1}(k)$}{абсолютный одноточечный критерий сходимости ландшафта функции потерь}
\nomenclature[33]{$\Delta_{2}(k)$}{среднеквадратичный критерий сходимости ландшафта функции потерь}
\nomenclature[34]{$\Delta_{2}^{(D)}(k)$}{среднеквадратичный критерий сходимости в подпространстве наибольшей кривизны размерности $D$}
\nomenclature[35]{$\mathcal{S}_{D} \subset \mathbb{R}^{N}$}{подпространство наибольшей кривизны размерности $D$}
\nomenclature[36]{$M_{G}(\mathcal{A})$}{локальная архитектурно-зависимая константа кривизны}
\nomenclature[37]{$\mathbf{X}_{m} = \left[ \mathbf{x}_1, \ldots, \mathbf{x}_m \right]^\top \in \mathbb{R}^{m \times n}$}{матрица объекты-признаки}
\nomenclature[38]{$\mathbf{y}_{m} = \left[ y_1, \ldots, y_m \right]^\top \in \mathbb{R}^{m}$}{вектор значений целевой переменной}
\nomenclature[39]{$L(\mathfrak{D}_m, \mathbf{w})$}{функция правдоподобия выборки}
\nomenclature[40]{$l(\mathfrak{D}_m, \mathbf{w})$}{логарифмическая функция правдоподобия выборки}
\nomenclature[41]{$\hat{\mathbf{w}}_{k} \in \mathbb{R}^{n}$}{оценка максимума правдоподобия по подвыборке $\mathfrak{D}_{k}$}
\nomenclature[42]{$\mathcal{I}(\mathbf{w}) \in \mathbb{R}^{n \times n}$}{информационная матрица Фишера}
\nomenclature[43]{$D_{\text{KL}}(p \| q)$}{дивергенция Кульбака~--"Лейблера}
\nomenclature[44]{$\text{s-score}$}{функция близости распределений $\text{s-score}$}

\printnomenclature[4.5cm] % Значение ширины столбца с обозначениями стоит подбирать вручную